\chapter{Discussion}
\label{chap:discussion}

%% -----------------------------------------------------------------------
\section{Answering the Research Questions}
\label{sec:disc:rq}

Three research questions (RQs) motivated this DDP.  This section
evaluates each against the experimental evidence.

\subsection{RQ1: Does joint inventory-transport optimisation outperform
single-axis optimisation?}

\textbf{Answer: Yes, decisively.}

Experiment E3 (joint optimisation) achieves a total cost of \$1,556,516
at a fill rate of 95.37\%.  This is:
\begin{itemize}
  \item \textbf{39.0\% below the analytical baseline (E0)}, which costs
        \$2,553,418 at only 86.15\% fill.
  \item \textbf{15.5\% below transport-only optimisation (E1b)}, which
        costs \$2,252,612 but achieves only 82.25\% fill.
  \item \textbf{42.4\% below inventory-only optimisation (E2)}, which
        costs \$2,700,388 even while achieving 97.33\% fill.
\end{itemize}

The mechanism is clear from the cost-shift analysis in
Section~\ref{sec:results:e3}: joint optimisation deliberately increases
holding cost by \$242k (+175\%) to unlock \$1.04M in transport savings
($-$55\%), achieving a 4.3:1 leverage ratio between holding and transport costs.

The reason single-axis optimisation cannot achieve this is structural.
Transport-only optimisation (E1b) cannot raise dispatch thresholds
without depleting retailer inventory, because base-stock levels are
anchored at lean analytical values.  Inventory-only optimisation (E2)
raises stock levels to achieve good fill rates but leaves dispatch
thresholds at zero, so every order is dispatched immediately as a small,
expensive partial load.  Only joint search discovers that raising stock
and raising dispatch thresholds are \emph{complementary} moves that must
be made simultaneously.

\subsection{RQ2: How does the joint-optimised configuration perform under
supply disruption?}
\label{sec:disc:rq2}

\textbf{This research question addresses a planned extension that has not
yet been conducted.}  The supply disruption sensitivity experiment
(§7.9 of the original research plan) will re-run E0 and E3 with a 14-day
supplier outage and compare the degradation in fill rate and cost between
the two configurations.  The hypothesis, consistent with the literature on
safety-stock buffering, is that the E3 configuration — which holds higher
inventory at the warehouse level — will degrade more gracefully than E0.
This experiment is reserved for future work and is described in the
limitations and future work sections below.

\subsection{RQ3: Does joint optimisation suppress the bullwhip effect?}

\textbf{Answer: Partially, but not fully.}

The bullwhip analysis in Section~\ref{sec:results:bullwhip} shows that
amplification ratios at the warehouse and supplier nodes decrease under
E3 relative to E0.  This is consistent with the mechanism: the high
dispatch thresholds on the supplier-to-warehouse lanes mean that the
warehouse accumulates demand before placing large consolidated orders,
which appear \emph{less} variable to the supplier on a per-day basis.

However, amplification is not eliminated.  The warehouse node still
aggregates orders from three retailers with heterogeneous demand patterns,
introducing variability that is not present in any individual retailer's
demand.  Eliminating bullwhip entirely would require either a
vendor-managed inventory arrangement (where the supplier directly observes
retailer demand) or demand signal sharing — both of which are outside the
scope of the current framework.

%% -----------------------------------------------------------------------
\section{Practical Recommendations}
\label{sec:disc:recommendations}

The experimental results translate into three actionable guidelines for
supply-chain practitioners:

\begin{enumerate}
  \item \textbf{Co-optimise inventory and dispatch when transport cost
        exceeds 50\% of total cost.}  In the E0 baseline, transport
        accounts for 74\% of total cost — a strong signal that transport
        is the primary lever.  Rule of thumb: if transport exceeds half
        the total cost, the joint optimisation approach is likely to
        return at least 20\% savings.

  \item \textbf{Accept a deliberate holding-cost increase to unlock
        consolidation.}  The 4.3:1 leverage ratio found in this study
        suggests that every additional dollar of safety stock invested
        in the warehouse returns approximately \$4.30 in transport
        savings.  This counter-intuitive trade-off is invisible to managers
        who optimise cost components independently.

  \item \textbf{Use the Pareto frontier to negotiate service-level
        agreements.}  The E4 results show a natural operating point at
        approximately 95\% fill / \$1.56M and a more expensive regime at
        98\% fill / \$1.68M, with sharply diminishing returns between the
        two.  When agreeing SLAs with customers, mapping this frontier
        allows the supplier to price each incremental service percentage
        point accurately.
\end{enumerate}

%% -----------------------------------------------------------------------
\section{Limitations}
\label{sec:disc:limitations}

The following limitations of the current framework should be considered
when interpreting the results.

\begin{description}
  \item[Deterministic lead times.] The simulator uses a fixed, deterministic
        transit time on each lane.  Real transport networks exhibit
        stochastic lead times driven by congestion, weather, and carrier
        variability.  Extending to stochastic lead times would increase
        the optimal safety-stock levels and may change the cost-service
        trade-off.

  \item[Single DC layer.] The 1N3 topology has one warehouse serving
        all three retailers.  More complex networks (two DCs, hub-and-spoke)
        may exhibit different consolidation dynamics.

  \item[CSV demand without forecasting.] Demand is read from historical
        M5 CSVs rather than from an online demand forecast.  In practice,
        policies are triggered by forecasts, not exact future demand, so
        safety stock should account for forecast error.

  \item[No lateral transshipment.] Retailers cannot transfer stock to
        each other.  Lateral transshipment is a common strategy for
        balancing inventory across retail locations and could substantially
        reduce shortage costs.

  \item[Single-SKU optimisation policy.] The main experiments use the
        base-stock policy for all nodes and SKUs.  The framework supports
        seven policies but the joint optimisation was not repeated for
        alternative policy families (e.g.\ $(k,m)$-cycle for all nodes).

  \item[Supply disruption experiment pending.] The analysis of how the
        optimised configuration degrades under supplier outage (RQ2)
        has not been conducted and is reserved for future work.
\end{description}
